%% conformal-antenna-diffopt.9. Sources: changelog.md, provenance.md, verify_evidence.py.

\documentclass{anvil-paper}
\usepackage{xurl}
\usepackage{array}

\expandafter\def\expandafter\UrlBreaks\expandafter{\UrlBreaks\do\-\do\_}
\Urlmuskip=0mu plus 1mu
\renewcommand{\topfraction}{0.92}
\renewcommand{\bottomfraction}{0.6}
\renewcommand{\textfraction}{0.07}
\renewcommand{\floatpagefraction}{0.85}
\setcounter{topnumber}{3}
\setcounter{totalnumber}{4}
\renewcommand{\bibfont}{\footnotesize}
\setlength{\bibsep}{0.5pt plus 0.5pt}

\newcommand{\anvilfig}[3][0.98]{%
  \IfFileExists{figures/#2}%
    {\includegraphics[width=#1\linewidth]{figures/#2}}%
    {\fbox{\parbox{#1\linewidth}{\centering\vspace{5em}%
      \texttt{figures/#2}\\[0.5em]%
      {\small placeholder: render via \texttt{figures/src/#3}}%
      \vspace{5em}}}}}
\newcommand{\PL}{P2/L1}

\title{Adjoint Shape Optimization of a Conformal 2.4\,GHz Patch Antenna:
EM Solves to Target Against Gradient-Free Search at Several Settings}
\author{Robb Walters\\
  {\normalsize Independent Researcher}}
\date{October 2026}

\begin{document}
\maketitle

\begin{abstract}
We design the copper of a probe-fed 2.4\,GHz patch antenna on a $1.6$\,mm
FR-4 board bent to a $100$\,mm radius with geode-fem, a finite-element
solver whose discrete shape adjoint reuses one factorization per frequency.
We measure cost in EM solves (one frequency, one sparse LU factorization) and
count the solves an optimizer spends until the worst $|S_{11}|$ over the band
first reaches a target level on the coarse optimization model. With
38 shape parameters the adjoint reaches $-9.5$\,dB in $35$ solves; the fastest
separable CMA-ES, full CMA-ES or NEWUOA run to reach that level needs $385$
($11\times$, or $8.3\times$ in geode-fem wall time, since an adjoint
evaluation costs $1.32$ forward ones), and the fastest to reach $-9$\,dB needs
$230$ against $30$ ($7.7\times$, full CMA-ES, one seed). At $-9.5$\,dB the
ratio is $2.4\times$, $2.9\times$ and $11\times$ with 4, 14 and 38 parameters.
L-BFGS with a finite-difference gradient follows the adjoint's iterates at
$39$ times its solves, behind CMA-ES, so the savings come from the adjoint.
Re-centred using Palace at third order and a rim-refined mesh, the design
covers $2.400$--$2.4835$\,GHz at $-10$\,dB on the three most refined models
computed, with $1.9$\,MHz of lower-edge margin on the finest; we claim
coverage on those models only. A larger air margin moves the resonance up by
$4.3$\,MHz and further refinement moves it up again, so we estimate that the
mesh-limit design misses the lower band edge by $1$ to $13$\,MHz.
\end{abstract}

\section{Introduction}
\label{sec:intro}

Gradient-based inverse design is routine in nanophotonics, where
differentiable solvers such as \textsc{ceviche} \citep{hughes2019forward},
the Meep adjoint solver \citep{hammond2022meepadjoint}, SPINS-B
\citep{su_spinsb} and FDTDX \citep{mahlau2026fdtdx} optimize a permittivity
density on a structured grid. Shape derivatives of scattering parameters for
finite-element models with ports are older
\citep{webb2000design,georgieva2002feasible,nikolova2004adjoint} and are
available in commercial tools
\citep{ansys_hfss_derivatives,comsol_sensitivity_method}. A gradient is worth
its implementation cost only if it buys something that gradient-free search
does not. With two nearly independent parameters, finite differences or a
scaling law are cheap. With many coupled parameters, finite differences cost
one solve per parameter and gradient-free search slows with dimension
\citep{rios2013derivative}, which is where an adjoint should pay.

This paper measures that on one antenna. A 2.4\,GHz probe-fed patch is bent
around a cylinder of radius $R_b = 100$\,mm. The board is a uniform
$1.6$\,mm FR-4 shell that is not designed; up to 38 smooth parameters move
the copper outline, the corners and the feed. geode-fem, a differentiable
finite-element solver with a discrete shape adjoint, optimizes the band
objective by L-BFGS. On the same problem, objective, start and volume guard
we run separable and full-covariance CMA-ES
\citep{hansen2001cmaes,ros2008separable} and NEWUOA
\citep{powell2006newuoa,zhang2023prima} over a range of settings, the same
L-BFGS with a finite-difference gradient, and, for two optimizers, the
independently developed solver Palace \citep{palace} as the forward model.

\paragraph{The cost measure.}
We measure an optimizer's cost in \emph{EM solves}. One EM solve is the
finite-element model at one frequency: one sparse LU factorization of the
$150{,}561$-unknown system of the coarse model and its back-substitutions.
The objective samples five frequencies, so one evaluation of it is five
solves; on one host it takes $72.1$\,s, and $95.1$\,s when the adjoint
gradient is computed with it, since the adjoint reuses each factorization for
one more back-substitution and then assembles the shape gradient. The quantity we compare is \emph{solves to
target}: the number of EM solves an optimizer has spent when its best design
so far first reaches a target level of the worst-of-band $|S_{11}|$, the
largest reflection over the band (Figure~\ref{fig:reading} shows how these
quantities are read off a response curve). Solves rather than seconds are the
measure because they do not depend on the host or on how a run is
parallelized; where the adjoint's higher cost per solve matters, we also
give the ratio in wall time.

Three results follow.
\begin{itemize}
  \item \textbf{The adjoint reaches each target in fewer EM solves, and the
  ratio is larger with more parameters} (Section~\ref{sec:baselines}). To
  $-9.5$\,dB the fastest gradient-free run needs $2.4\times$, $2.9\times$ and
  $11\times$ the adjoint's solves with 4, 14 and 38 parameters, and to
  $-9$\,dB $1.1\times$, $3.3\times$ and $7.7\times$. The adjoint's count to
  $-9.5$\,dB falls ($85$, $40$, $35$) while the best gradient-free count does
  not follow a trend ($205$, $115$, $385$); three sizes of one problem show
  the direction, not a rate. With 4 parameters three gradient-free runs end
  lower than the adjoint. A finite-difference gradient on the same L-BFGS path
  costs $39$ times the adjoint's solves, so the savings are the adjoint's, not
  the gradient's as such; given enough solves, gradient-free search reaches
  the same level.
  \item \textbf{The forward model's error sits at the patch rim}
  (Section~\ref{sec:accuracy}), the thin region along the copper's edge
  (Figure~\ref{fig:geometry}d). Palace on geode-fem's own meshes agrees with
  geode-fem at first and second order, and at third order with geode-fem's
  finest rim-refined mesh, and a port-free eigenmode places the matched dip
  (the frequency of minimum $|S_{11}|$) $16$\,MHz above the cavity resonance.
  \item \textbf{A design re-centred on the refined model covers the ISM band
  on the most refined models computed, and is estimated to miss it in the
  mesh limit} (Section~\ref{sec:results}). Measured: on the three most refined
  models its $-10$\,dB band covers $2.400$--$2.4835$\,GHz, with $1.9$ to
  $6.6$\,MHz to spare at $2.400$\,GHz. Estimated: a larger air margin and
  further refinement move the band up, and we estimate that the mesh-limit
  design misses $2.400$\,GHz by $1$ to $13$\,MHz. We therefore claim ISM
  coverage for those discrete models, not for the mesh-limit model or a built
  antenna. The probe has a realizable width ($1.6$\,mm, modelled as a lumped
  port).
\end{itemize}

\paragraph{Evidence tiers.}
We keep three tiers apart. A \emph{bounded} result is a proven two-sided
bound; no result in this paper is of that kind. A \emph{model-dependent}
result is a value computed for a stated discrete model (mesh, element order,
absorber and port model). An \emph{estimated} result is an extrapolation to
the mesh limit, given with its band and never called a converged value.

\paragraph{Names.}
The paper refers to parts of the antenna, features of its response and
stages of the design by name. Table~\ref{tab:terms} collects them, each with
the figure or table that shows it.

\begin{table}[t]
  \centering
  \caption{Names used in the text, with where each is shown.}
  \label{tab:terms}
  \footnotesize
  \begin{tabular}{@{}>{\raggedright\arraybackslash}p{0.17\linewidth}>{\raggedright\arraybackslash}p{0.62\linewidth}>{\raggedright\arraybackslash}p{0.15\linewidth}@{}}
    \toprule
    Name & Meaning & Shown in \\
    \midrule
    Patch, ground & The copper sheet on top of the board (the only part designed) and the copper on its underside (fixed) & Fig.~\ref{fig:geometry}b, c \\
    Board, substrate, slab & The $1.6$\,mm FR-4 shell between patch and ground, bent to $R_b = 100$\,mm; the slab is the substrate under the patch & Fig.~\ref{fig:geometry}b \\
    Patch edges $\pm y$, $\pm s$ & The four sides of the patch. The $\pm y$ edges end the resonant length and radiate; the $\pm s$ edges run across the bend and do not & Fig.~\ref{fig:geometry}b, c, e \\
    Rim & The region within $0.5$\,mm of the patch edge, where the field is singular and the mesh is refined & Fig.~\ref{fig:geometry}d; Table~\ref{tab:x1} \\
    Probe, feed, port & The $1.6$\,mm-wide connection from ground to patch near the $+y$ edge, modelled as a lumped port; the inset is its distance from that edge & Fig.~\ref{fig:geometry}c, e \\
    Air margin, UPML & The air gap between the board and the absorbing layer (UPML) that stands in for open space & Fig.~\ref{fig:geometry}a \\
    Sagitta & How far the patch edge lies below the patch's crown because of the bend & Fig.~\ref{fig:geometry}b \\
    Chart $(s, y)$ & Coordinates on the board: $s$ is arc length across the bend, $y$ runs along the bend axis & Fig.~\ref{fig:geometry}e, f \\
    Mesh morphing, volume guard & Moving mesh nodes instead of remeshing; the guard rejects a step that would nearly invert an element & Section~\ref{sec:freedom}, \ref{sec:objective} \\
    Dip & The frequency of minimum $|S_{11}|$, where the antenna is best matched (the resonance as seen at the port) & Fig.~\ref{fig:reading} \\
    $-10$\,dB band, band edges, margins & The range where $|S_{11}| \le -10$\,dB; its ends (the $-10$\,dB crossings); their distances outside the ISM band edges & Fig.~\ref{fig:reading} \\
    Worst-of-band $|S_{11}|$, level & The largest $|S_{11}|$ over the band points; a level is a target value for it & Fig.~\ref{fig:reading}, \ref{fig:optimizers} \\
    EM solve, solves to target & One frequency of the model (one LU factorization); solves spent until a target level is first reached & Section~\ref{sec:intro}; Fig.~\ref{fig:optimizers} \\
    L1, \PL, refinement levels & The coarse optimization mesh; L1 at element order $p = 2$; L1 with a finer slab, rim or both & Table~\ref{tab:x1} \\
    Space mapping, $\delta$ & Optimizing on L1 for a band shifted down by $\delta$, the fitted coarse-to-fine frequency offset & Section~\ref{sec:procedure} \\
    Designs A, B, C & The ends of the design chain; C is the final design & Table~\ref{tab:chain}; Fig.~\ref{fig:geometry}f \\
    \bottomrule
  \end{tabular}
\end{table}

\section{Related Work}
\label{sec:related}

\paragraph{Finite-element S-parameter sensitivities with ports.}
Design sensitivities of the scattering matrix for tetrahedral edge elements
were derived by \citet{webb2000design}, extended to the frequency response
by \citet{webb2002design} and used in gradient optimization of microwave
devices by \citet{nair2003optimization}. \citet{georgieva2002feasible} give
adjoint S-parameter sensitivities that need only the solver's field
solutions, so that commercial solvers can supply them; adjoint and
self-adjoint formulations are reviewed by \citet{nikolova2004adjoint} and
\citet{nikolova2006network}, and continuum shape sensitivity for
high-frequency devices is given by \citet{choi2011generalized}. HFSS computes
S-parameter derivatives with respect to geometric parameters by reusing the
LU factors \citep{ansys_hfss_derivatives,vardapetyan2008sensitivity}, and
COMSOL documents adjoint shape optimization of waveguide devices on a
deformed mesh \citep{comsol_sensitivity_method,jensen2022comsolshape}. For
antennas, \citet{ghassemi2013antenna} evolve a microstrip antenna with an FEM
adjoint sensitivity and a small set of control vertices.
\citet{toivanen2009electromagnetic} obtain antenna shape sensitivities from a
method-of-moments solver by automatic differentiation, and
\citet{toivanen2010gradient} optimize ultra-wideband antennas whose outlines
are parametrized by splines, a close prior combination to our spline
outline and gradient descent. \citet{jacobsson2010gradient} apply
gradient-based shape optimization to conformal array antennas.

\paragraph{Gradient-free and surrogate-assisted antenna optimization.}
Evolutionary search is the common gradient-free route in antenna design, and
surrogate models cut its cost: \citet{liu2014efficient} combine differential
evolution with Gaussian-process surrogates, \citet{zhou2021trust} use
trust-region parallel Bayesian optimization, and \citet{akinsolu2020machine}
review machine-learning-assisted methods. The closest precedent for our
measurement is aerodynamic: \citet{zingg2008comparative} compare a genetic
algorithm with an adjoint method on aerodynamic shape problems and find that the
genetic algorithm needs $5$ to $200$ times as many function evaluations. We
did not run a surrogate-assisted or Bayesian optimizer.

\paragraph{Variable-fidelity antenna optimization and space mapping.}
Space mapping aligns a cheap coarse model with an accurate fine model and
optimizes the coarse one \citep{bandler1994space,bandler2004space,koziel2006framework}.
In antenna design the coarse model is often the same electromagnetic solver
on a coarse discretization, corrected by frequency scaling or by
response-shape correction \citep{koziel2010sprp,koziel2014antenna}, and
adjoint sensitivities have been used inside trust-region space mapping
\citep{koziel2013reliable}. Our correction is the simplest member of this
family: a frequency shift fitted between two discretizations of the same
solver.

\paragraph{Automated shape differentiation and structured grids.}
Shape derivatives of finite-element models can be generated automatically in
UFL and Firedrake \citep{ham2019automated,dokken2020,paganini2021fireshape}
and in NGSolve \citep{gangl2021fully}. \citet{wang2011} give a discrete
Maxwell shape adjoint on unstructured meshes in two dimensions, in the time
domain and without ports. On structured grids, conductors are designed by
density-based topology optimization
\citep{erentok2011topology,hassan2014topology,hassan2014patch} and with FDTD
shape sensitivities \citep{nikolova2004sensitivity,nikolova2006scattering}.
A plain Yee grid staircases a curved conductor
\citep{cangellaris1991analysis,haggblad2014accuracy}, which locally
conformal FDTD avoids
\citep{dey1997locally,mittra2000conformal,zagorodnov2003uniformly,benkler2006new}.
Subpixel-smoothed projection gives the Meep density adjoint
subpixel-accurate boundaries \citep{hammond2025unifying,romano2026ssp2},
though its material-grid adjoint excludes metals \citep{meepadjointdocs};
Tidy3D offers autograd shape gradients and conformal PEC averaging
\citep{tidy3dautograd}; shape parameters can be mapped to a fixed grid
\citep{hooten2025,liu2024imagerep}; and boundary-integral methods carry
recent PEC shape-sensitivity work
\citep{takahashi2025tdbie,arens2026tubular,balouchev2024balun}.

\paragraph{Conformal and cylindrical patch antennas.}
The cylindrical-rectangular patch has cavity-model analyses
\citep{krowne1983cylindrical,luk1989cylindrical}, and nonplanar microstrip
antennas and conformal arrays are treated in the monographs of
\citet{wong1999nonplanar} and \citet{josefsson2006conformal}. We use a
full-wave model and do not compare with these analyses; the curvature here
is mild (Section~\ref{sec:accuracy}).

\paragraph{Position of this work.}
What this paper adds is a measurement on an open-radiator problem with a
moving feed: the solve cost of a finite-element shape adjoint against
gradient-free optimizers and a finite-difference gradient at 4, 14 and 38
parameters. We do not claim a first EM shape adjoint, and we make no
performance comparison with structured-grid methods.

\section{Method}
\label{sec:method}

\begin{figure}[tp]
  \centering
  \anvilfig{fig1_geometry.pdf}{fig1_geometry.py}
  \caption{The fixture and its named parts (Table~\ref{tab:terms}). (a)~The
  board inside the air box and the UPML, to scale. (b)~Across the bend: the
  slab is the substrate under the patch (dotted); circles mark the rim.
  (c)~Through the feed: the probe at its initial inset. (d)~The rim, the
  region within $0.5$\,mm of the patch edge whose mesh Table~\ref{tab:x1}
  refines. (e)~What the 38 parameters move on the initial patch: each small
  arrow is one B-spline coefficient (positions schematic), the diagonal
  arrows the corner bumps, plus the probe's inset and width $w_p$. (f)~The
  initial patch (dashed) and design~C (filled). The board is the same in
  every design. Drawn from \texttt{baseline\_2g4\_L1.toml} and
  \texttt{recentre/geometry/rbG\_L1\_copper\_38\_sm57\_w16.json}.}
  \label{fig:geometry}
\end{figure}

\subsection{The open-radiator forward model}
\label{sec:pencil}

The forward problem is the driven frequency-domain curl-curl equation,
discretized with N\'ed\'elec edge elements of order $p = 1$ or $p = 2$ on an
unstructured tetrahedral mesh. Lengths are in millimetres and the solver
works with $\omega = k_0 = 2\pi f/c$ in rad/mm. The assembled system is the
complex-symmetric pencil
\begin{equation}
  A(\omega)\,\mathbf{x}
  = \Big[\,K(\nu) - \omega^2 M(\varepsilon) + \tfrac{j\omega}{Z_s}\,S_p\,\Big]\,\mathbf{x}
  = \mathbf{b},
  \label{eq:pencil}
\end{equation}
with curl-curl stiffness $K$, mass matrix $M$ and lumped-port surface term
$S_p$. The radiating boundary is a box UPML of thickness $t = 25$\,mm:
per-element tensor weights $\Lambda$ enter $K$ (as $\nu = \Lambda^{-1}$) and
$M$ (as $\varepsilon = \varepsilon_r \Lambda$), with the stretch
$\varsigma = 1 - j\sigma_0 (d/t)^2/\omega$ per axis at depth $d$ and
$\sigma_0 = 1$. The FR-4 has $\varepsilon_r = 4.4$ and $\tan\delta = 0.02$,
and the conductors are perfect electric conductors (PEC). The port is a
rectangle of width $w_p$ along $y$ and length $l_p = 1.6$\,mm across the
substrate, with sheet impedance $Z_s = R_0\,w_p/l_p$ and $R_0 = 50\,\Omega$.
The reflection coefficient follows from the port voltage $V$:
\begin{equation}
  S_{11} = \frac{Z - R_0}{Z + R_0},\qquad Z = \frac{V}{I},\qquad
  I = \frac{2V_{\mathrm{inc}} - V}{R_0}.
  \label{eq:s11}
\end{equation}
Every evaluation checks passivity, $|S_{11}| \le 1$.

\subsection{Shape adjoint with one factorization per frequency}
\label{sec:adjoint}

The design moves mesh node coordinates $X$, so the gradient needs
$\partial K/\partial X$, $\partial M/\partial X$, $\partial S_p/\partial X$
and $\partial \mathbf{b}/\partial X$. geode-fem evaluates its closed-form
element kernels in first-order dual arithmetic inside a tensor stack
\citep{burn}: seeding the tangent of one
node coordinate returns the exact partial derivatives of every local matrix
entry. The UPML weights are held fixed under node motion. Because $A(\omega)$
is complex symmetric, the adjoint system uses the same operator, and one LU
factorization per frequency serves the forward and the adjoint solve. The
nodal gradient is $\lambda^{\top}(\partial A/\partial X)\,\mathbf{x}$ plus the
explicit dependence of the port drive and of the port-voltage functional on
the geometry. When a parameter moves the feed, the port faces move and the
port is rebuilt on every moved mesh. A unit test of the moving-feed gradient
on a curved mesh, with a source that varies across the port so that the
explicit term matters, matches central finite differences to below
$10^{-8}$; without the explicit term it misses by a relative error of about 3.

\subsection{Copper design freedom on a fixed curved board}
\label{sec:freedom}

\paragraph{Mesh morphing.}
The tetrahedral mesh conforms to the copper: the patch outline is made of
mesh edges. A design is not remeshed. It moves the mesh nodes instead, which is called mesh
morphing: nodes on and under the copper move with it, the air nodes around
them follow smoothly, and the absorber and the outer boundary stay fixed.
Because the node positions are smooth functions of the design parameters,
the assembled system, and with it $|S_{11}|$, is differentiable in them, and
the adjoint of Section~\ref{sec:adjoint} gives that derivative. Morphing has
a limit: where a moving edge approaches a fixed feature, the elements between
them are squeezed, and a tetrahedron squeezed to zero volume turns inside out
(inverts) and invalidates the discretization. The volume guard of
Section~\ref{sec:objective} keeps every design short of that.

\paragraph{The morph.}
The board is not a design variable. Each substrate node has chart
coordinates $(s, y)$, with $s = R_b\,\mathrm{atan2}(x, z - z_{\mathrm{axis}})$
the mid-plane arc length and $y$ the bend axis (Figure~\ref{fig:geometry}e). A design applies a chart
displacement $(\Delta s, \Delta y)$ and moves the node to angle
$\phi_0 + \Delta s/R_b$ and height $y_0 + \Delta y$ at its own radius, so
the FR-4 stays a uniform $1.6$\,mm shell. Every run records the largest
change of any substrate node's radius ($\le 1.4\times10^{-14}$\,mm). Air
nodes follow a harmonic extension of the motion; the UPML and the outer
boundary are fixed. The design gradient chains the nodal adjoint gradient
through the exact Jacobian of this map.

\paragraph{The parameters.}
The 38 parameters (Figure~\ref{fig:geometry}e), each per unit coefficient, are $8$ clamped cubic
B-splines per patch edge ($\pm y$, $\pm s$), each moving its edge outward in
the surface by $1$\,mm, ramped in from the feed band and decaying over the
$12$\,mm substrate margin (32); one bump per corner, moving it $1$\,mm
outward along the bisector with support $1 - \mathrm{smoothstep}(d/6\,\mathrm{mm})$
(4); and the probe's inset from the $+y$ edge and its width $w_p$, each
$1$\,mm per unit (2). They
are smooth functions of the chart, so the same coefficient vector defines
the same copper on every mesh of the family, which lets a design be carried
between discretizations. The basis cannot open a hole in the patch. For the
scaling study two smaller bases are used: 14 parameters (two linear outline
splines per edge, the corners and the feed) and 4 (patch width and length,
probe inset and probe width); they are separate bases, not subsets of the
38.

\subsection{Objective, optimizer and guards}
\label{sec:objective}

The objective is a soft maximum of $|S_{11}|^2$ over $n_f = 5$ band
frequencies, equally spaced from $2.400$ to $2.4835$\,GHz,
\begin{equation}
  G(\mathbf{c}) = \frac{1}{\beta}\ln \sum_{f} \exp\!\big(\beta\,|S_{11}(f;\mathbf{c})|^2\big),
  \qquad \beta = 25,
  \label{eq:objective}
\end{equation}
over the coefficient vector $\mathbf{c}$. $G$ exceeds the largest
$|S_{11}|^2$ by at most $\ln 5/\beta = 0.064$, and every level reported in
this paper is the worst of the band points, not $G$
(Figure~\ref{fig:reading}). L-BFGS (memory 8) with
Armijo backtracking and a cap of 1 coefficient unit per step minimizes $G$.

\begin{figure}[t]
  \centering
  \anvilfig{fig2_reading.pdf}{fig2_reading.py}
  \caption{How a return-loss curve is read, and the names the text uses for
  its features. $|S_{11}|$ is the fraction of incident power reflected at the
  port, in dB ($-10$\,dB: $10\%$ reflected). Shown are design~C
  (Section~\ref{sec:procedure}) on Palace at $p = 3$ on L1 (solid) and the
  initial patch on L1 at $p = 1$ (dashed). The dip ($\blacktriangledown$) is the
  frequency of minimum $|S_{11}|$. The $-10$\,dB band runs between the two
  $-10$\,dB crossings, the band edges; the margins are how far the band edges
  lie outside the ISM band (shaded). The worst-of-band $|S_{11}|$ is the
  largest value over the in-band points (here the nine verification points;
  the objective uses five), and a level is a target value for it. The initial
  patch's dip lies $89$\,MHz below the band, and its worst-of-band value is
  near $-1.5$\,dB. Data: \texttt{baseline\_2g4\_L1.toml},
  \texttt{recentre/palace/sm57\_p3/port-S.csv},
  \texttt{recentre/recentre\_summary.toml}.}
  \label{fig:reading}
\end{figure}

\paragraph{The volume guard and how a descent ends.}
Each trial step of the line search is first morphed and checked: if it would
shrink any tetrahedron's signed volume below $0.25$ of its volume in the
unmorphed mesh, the step is rejected without a solve and halved. A volume
ratio of $0.25$ leaves a margin before inversion at $0$, so the guard bounds
how far the mesh is allowed to be morphed. A descent \emph{ends on the guard}
when no trial step both passes the guard and lowers $G$ enough within 25
halvings while some trial failed the guard, or when it reaches a plateau with
the most compressed element within $5\%$ of the guard. It \emph{ends on a
plateau} when four accepted steps in a row each lower $G$ by less than
$5\times10^{-4}$ of its value. Every adjoint descent reported here ends on the
guard, except one intermediate run of the space-mapping chain
(Section~\ref{sec:procedure}), which ended on a plateau. A descent that ends
on the guard has stopped where this mesh can be morphed no further along its
path, not necessarily at a stationary point of the design problem; the guard's
site and what would relieve it are given in Sections~\ref{sec:procedure}
and~\ref{sec:results}.

Before and after each run the pipeline gradient is compared with a central
finite difference along the steepest-descent direction and along random
directions. The worst relative error over the runs reported here is
$1.5\times10^{-5}$ (tolerance $5\times10^{-3}$), and a per-group check on the
two coarsest meshes passes to $1.9\times10^{-7}$. The adjoint is implemented
for $p = 1$ only, so optimization runs at $p = 1$ and finer models enter
through verification.

\section{Fixture and Forward-Model Accuracy}
\label{sec:accuracy}

\paragraph{Antenna.}
The initial antenna is a $38\times29$\,mm PEC patch (width along $s$,
resonant length along $y$) on $1.6$\,mm FR-4 with a $12$\,mm substrate and
ground margin, fed by a $1.6$\,mm probe, modelled as a uniform lumped port,
inset $7.0$\,mm from the $+y$ edge (Figure~\ref{fig:geometry}c). The bend axis
is parallel to $y$, so the curvature runs along the non-radiating width; the
sagitta at the patch edge is $1.8$\,mm (Figure~\ref{fig:geometry}b). A
$40$\,mm air margin separates the board from the UPML
(Figure~\ref{fig:geometry}a). The target is $|S_{11}| \le -10$\,dB over the
$2.400$--$2.4835$\,GHz ISM band.

\paragraph{The coarse mesh.}
The optimization mesh L1, level 1 of the fixture's uniformly refined mesh
family L0 to L3, has $150{,}561$ edges, a median of $2.0$ element
layers through the slab and a mean edge length of
$0.78$\,mm at the patch rim. On L1 at $p = 1$ the initial patch's
$|S_{11}|$ dip (Figure~\ref{fig:reading}) is at $2.3118$\,GHz; at $p = 2$ on L1 (\PL) it is at
$2.3566$\,GHz, $45$\,MHz higher. A direct LU solve at \PL{} ($804{,}460$
unknowns) takes $132$\,GiB with geode-fem's direct solver, so refining
further at $p = 2$ is not practical with this solver.

\begin{table}[t]
  \centering
  \caption{Where the error sits: the initial patch's $|S_{11}|$ dip under
  refinement of L1, geode-fem at $p = 1$, with \PL{} for reference. Each row
  is one refinement level. ``Slab'' is the element size in the substrate
  under the patch (Figure~\ref{fig:geometry}b), ``rim'' the mean edge length
  within $0.5$\,mm of the patch edge (Figure~\ref{fig:geometry}d), ``layers''
  the median number of elements through the slab. Model-dependent values.}
  \label{tab:x1}
  \small
  \begin{tabular}{@{}l c c c r c r@{}}
    \toprule
    Mesh & Slab (mm) & Rim (mm) & Layers & DOFs & Dip (GHz) & vs L1 (MHz) \\
    \midrule
    L1 & $0.80$ & $0.78$ & $2.00$ & $150{,}561$ & $2.3118$ & $0$ \\
    slab $0.533$ & $0.533$ & $0.52$ & $2.54$ & $234{,}691$ & $2.3273$ & $+15.5$ \\
    slab $0.40$ & $0.40$ & $0.39$ & $3.44$ & $414{,}837$ & $2.3404$ & $+28.6$ \\
    slab $0.30$ & $0.30$ & $0.29$ & $4.19$ & $828{,}775$ & $2.3480$ & $+36.2$ \\
    rim $0.40$ & $0.80$ & $0.39$ & $1.99$ & $317{,}424$ & $2.3420$ & $+30.2$ \\
    rim $0.20$ & $0.80$ & $0.20$ & $1.99$ & $430{,}808$ & $2.3569$ & $+45.1$ \\
    slab $0.30$ + rim $0.10$ & $0.30$ & $0.10$ & $4.19$ & $1{,}738{,}557$ & $2.3659$ & $+54.1$ \\
    \midrule
    L1 at $p = 2$ (\PL) & $0.80$ & $0.78$ & $2.00$ & $804{,}460$ & $2.3566$ & $+44.8$ \\
    \bottomrule
  \end{tabular}
\end{table}

\paragraph{The error sits at the rim.}
Table~\ref{tab:x1} refines L1 in two ways. Refining the whole slab under the
patch to $0.40$\,mm gives $3.4$ layers and, as a side effect, a $0.39$\,mm
rim. Refining only within $0.5$\,mm of the patch edge to the same rim size,
with the 2-layer interior unchanged, lands within $1.6$\,MHz of it
($2.3404$ against $2.3420$\,GHz). The interior layer count does not decide
the dip; the rim does, as expected where the field is singular at a
conductor edge \citep{meixner1972behavior}. A $0.2$\,mm rim at $p = 1$
($431$k unknowns) matches \PL{} ($805$k) to $0.3$\,MHz. Over all seven
$p = 1$ refinement levels of Table~\ref{tab:x1} the dip follows $f = f_\infty - C\,h_{\mathrm{rim}}^{\,0.80}$
with a $1.3$\,MHz rms residual (Figure~\ref{fig:accuracy}b). The exponent is
a fit to the rim edge length on graded meshes, not a proven
asymptotic order. The combined level, which halves the rim to $0.1$\,mm and
refines the slab, moves the dip $9.0$\,MHz.

\begin{figure}[t]
  \centering
  \anvilfig{fig3_accuracy.pdf}{fig3_accuracy.py}
  \caption{The initial patch's $|S_{11}|$ dip under refinement. (a) Against
  the number of unknowns: geode-fem at $p = 1$ on L1 and its slab, rim and
  combined refinements, geode-fem at $p = 2$ on L1, and Palace at
  $p = 1, 2, 3$ on the same L1 mesh. (b) Against the mean edge length at the
  patch rim, with the fit $f_\infty - C h^{0.80}$ (dashed), its limit and the
  order-1 Richardson step on the $0.39$/$0.20$\,mm rim pair. The shaded band,
  $2.366$--$2.380$\,GHz, runs from the finest computed value to the fit's
  limit; it is an estimate, not a bound. Data:
  \texttt{v6/v6\_summary.toml} \texttt{[x1]}.}
  \label{fig:accuracy}
\end{figure}

\paragraph{An independent solver on the same meshes.}
Palace \citep{palace} solved the same problem on geode-fem's own meshes,
written out with their physical tags, with the same materials, PEC surfaces
and uniform $50\,\Omega$ lumped port. Its absorber differs (air in the UPML shell, a first-order absorbing outer
wall); its driven sweeps use an adaptive reduced-order model with tolerance
$10^{-3}$ on a $1$\,MHz grid, and GMRES to $10^{-8}$. On L1, Palace and
geode-fem agree on the dip to $1.1$\,MHz at $p = 1$ ($2.3130$ against
$2.3118$\,GHz) and $0.4$\,MHz at $p = 2$ ($2.3570$ against $2.3566$\,GHz);
on the strip-fed design of Section~\ref{sec:procedure} the \PL{} band edges
agree to $0.7$\,MHz (lower) and $0.2$\,MHz (upper). Palace at $p = 3$ on L1
gives $2.3660$\,GHz, $8.9$\,MHz above its $p = 2$ value, and geode-fem's
finest refinement level (slab $0.30$ plus rim $0.10$\,mm) gives $2.3659$\,GHz
(Figure~\ref{fig:accuracy}a): two different refinements of the same model,
order in Palace and rim size in geode-fem, land within the $1$\,MHz sweep
step. Sharing mesh, port model and geometry, this checks the discretization,
not the model.

\paragraph{An FDTD cross-check.}
A Meep \citep{oskooi2010meep} model of the initial antenna checks the
resonance with an independent method. It puts zero-thickness PEC sheets on
the Yee grid, keeps the $1.6$\,mm gap in every column (4 cells at $0.4$\,mm, 8
at $0.2$\,mm), uses a port of exactly $50.0\,\Omega$ and models the FR-4 loss
as a conductivity ($\tan\delta = 0.0195$--$0.0205$ over $2.30$--$2.42$\,GHz);
an earlier Meep model with one-cell metal shells, which give the normal and
tangential field different gaps on the curved board, placed the bent dip
about $250$\,MHz low. The flat board dips at $2.3683$\,GHz, $10.3$\,MHz above
geode-fem's finest flat model ($p = 2$ on L1, itself about $9$\,MHz below
$p = 3$ for the bent patch), and the bend moves the dip by $-1.6$\,MHz,
against geode-fem's $-1.0$ to $-1.6$\,MHz. The bent dip is at $2.3667$\,GHz
with $0.4$\,mm cells and $2.3611$\,GHz with $0.2$\,mm cells
($\mathrm{Re}\,Z$ $44.8$ and $47.3\,\Omega$), against $2.3660$\,GHz for Palace
$p = 3$ and for geode-fem's finest refinement level. Meep's dips are parabola vertices
through samples $10$\,MHz apart, so differences of about $1$\,MHz are below
their resolution. The two Meep grids also differ in patch length ($28.8$
against $29.0$\,mm), width ($37.2$ against $37.6$\,mm, nominal $38$) and probe
inset ($7.0$ against $6.9$\,mm), so their $-5.6$\,MHz change is a trend, not a
mesh-limit value. The refined models of the three codes place the bent dip
between $2.3611$ and $2.3667$\,GHz.

\paragraph{The mesh-limit resonance.}
Every refinement computed so far raises the dip. If that continues, the
finest computed value, $2.366$\,GHz, is a lower estimate of the mesh-limit
dip; this is an inference from the trend, not a proof. The rim fit puts the limit at $2.380$\,GHz and the order-1
Richardson step on the rim pair at $2.372$\,GHz; since the fitted exponent
is below 1, the order-1 step is the lower of the two if the fit holds. We
therefore estimate the initial patch's mesh-limit dip at
$2.366$--$2.380$\,GHz, $0$ to $14$\,MHz above the finest computed models.

\paragraph{Port-free resonance.}
Palace's eigenmode solver, on the same L1 mesh at $p = 2$ with the port gap
left open (no feed), finds the patch mode at $2.3408$\,GHz with $Q = 22.2$
for the initial patch, and at $2.3991$\,GHz with $Q = 19.3$ for the final
design (7 modes, backward error below $5\times10^{-9}$), identified by its
substrate energy fraction ($0.80$--$0.83$, every other mode at most
$0.013$). The matched dip on the same discrete model lies above it, by
$16.3$\,MHz for the initial patch and $30.9$\,MHz for the design. In this
model the port is a resistive sheet: loading the gap with it raises the mode
by $12.2$ and $23.1$\,MHz and halves its $Q$ ($11.5$ and $8.6$), and the driven
dip sits a further $4.1$ and $7.7$\,MHz above the loaded mode. A series
inductance, as in a coaxial probe, would move the match the same way.
The Hammerstad transmission-line estimate for the initial patch,
$2.433$\,GHz \citep{hammerstad1975equations}, lies $92.5$\,MHz above the
port-free resonance at $p = 2$ and $76$\,MHz above the matched dip. The
discretization accounts for $9$ to $22$\,MHz of that and the feed moves the
dip toward the closed form, so most of the gap is in the cavity model; the
finite ground and the closed form's own accuracy are untested candidates.

\paragraph{Absorber, air margin and curvature.}
Over a ladder of UPML conductivity and thickness, $|S_{11}|$ changes by at
most $0.021$. Re-meshing the final design with a $60$ or $80$\,mm air margin
instead of $40$\,mm, at $p = 1$ on L1, moves its dip by $4.3$\,MHz at both
($4.34$ and $4.31$\,MHz), so the resonance has levelled off by $60$\,mm. The
band edges scatter about that between the three air meshes: the lower edge
rises $3.0$ and $4.3$\,MHz, the upper $4.2$ and $2.9$\,MHz, and the $-10$\,dB
bandwidth is $91.0$, $92.3$ and $89.7$\,MHz. The air-margin term is therefore
about $4.3$\,MHz with about $1.3$\,MHz of re-meshing scatter in the edges,
smaller than the refinement shifts but of the size of the design's margins
(Section~\ref{sec:results}).
Design~A's copper (the strip-fed design of Section~\ref{sec:procedure}, Table~\ref{tab:chain}) on
a flat board (bend radius $10^6$\,mm, the same mesh node for node)
resonates $0.3$\,MHz from the bent one at $p = 1$, and a complex
frequency-shift fit at \PL{} finds $0.28$\,MHz. The bend does not detune
that design, but it changes its match: the bent board has a $93.3$\,MHz
$p = 1$ $-10$\,dB band against $82.8$\,MHz flat, and a worst in-band
$|S_{11}|$ at \PL{} of $-10.78$ against $-9.93$\,dB. Design~C was not run
flat, and we did not investigate the mechanism.

\section{The Adjoint Against Gradient-Free Search}
\label{sec:baselines}

\paragraph{Setup.}
Every optimizer solves one problem: the objective (\ref{eq:objective}) on
the unshifted band, the $0.25$ volume guard and the L1 fixture at $p = 1$,
from $\mathbf{c} = 0$ ($-1.47$\,dB). The finite-difference control runs the
adjoint's L-BFGS, line search, step cap and guard with a one-sided
difference gradient (step $10^{-5}$), one band evaluation per parameter.
Separable CMA-ES is geode-fem's own implementation of
\citet{ros2008separable}, not cross-checked against a reference
implementation on a test function; full-covariance CMA-ES is pycma
\citep{pycma} with its defaults. Both use population
$\lambda = 4 + \lfloor 3\ln n\rfloor$ (14 at $n = 38$) and rank a sample
that fails the guard after every feasible sample, by its violation, without
a solve. NEWUOA runs in PRIMA with trust radii $0.5$ to $10^{-3}$ and $n + 2$
or $2n + 1$ interpolation points, and a design that fails the guard receives
a penalty above any feasible objective, without a solve. For Palace,
geode-fem applies the morph and the guard and writes the moved mesh, and
Palace solves the five band frequencies at $p = 1$ with GMRES on 16 MPI
ranks, a configuration not tuned for speed. The Palace CMA-ES runs reuse the
three geode-fem seeds at $\sigma_0 = 0.5$; with the same seed the two solvers
make the same evaluations to about evaluation 70, so those six runs are
three trajectories, each under two forward models, which isolates the
forward solver; seed 2 also reaches $-9.5$\,dB at the same count under both.
By ``several settings'' we mean two or three values of each optimizer's
main setting (Table~\ref{tab:baselines}); the best CMA-ES setting is the
smallest $\sigma_0$ tried. As defined in Section~\ref{sec:intro}, one EM solve is one band frequency with one
matrix factorization, and a level is a target for the worst-of-band $|S_{11}|$. The adjoint
solve reuses the forward factorization, and the adjoint's start, which
yields its first gradient, is counted; the run's recorded $190$ driven
solves include $65$ that are not counted: $35$ for the three gradient checks
and a fresh evaluation of the reported design, and $30$ for the verification
sweeps ($9$ in-band and $21$ wide points). An
adjoint band evaluation takes $95.1$\,s against $72.1$\,s for a forward one on
the same host ($1.32\times$), so in geode-fem wall time the $11\times$ at
$-9.5$\,dB below is about $8.3\times$ ($5{,}552$ against $666$\,s). The
finite-difference control is
charged $5$ solves per guard-passing trial plus $190$ per gradient. The
gradient-free runs' separate evaluation at $\mathbf{c} = 0$ is not counted,
except as NEWUOA's own first evaluation, and a sample rejected by the guard
costs no solve. Each level is the worst of the five band points in that
run's own forward model, and a level a run did not reach is censored at the
solves it used.
\begin{table}[t]
  \centering
  \caption{EM solves to reach fixed worst-of-band levels (target values of the
  worst $|S_{11}|$ over the band) with 38 parameters on the L1
  optimization model ($p = 1$), from $\mathbf{c} = 0$. Slashes separate
  seeds; the two $\sigma_0 = 0.5$ separable CMA-ES rows use the same three
  seeds in the same order and coincide to about evaluation 70. ``Used'' is
  the solves a run used, where a censored entry (--) stops. Budgets: 190
  evaluations; 114 for the first $\sigma_0 = 0.5$ seed on geode-fem, 1000
  for the long run, $2{,}500$ solves for finite differences, checked after
  each gradient. ``Best'' is over the counted evaluations, which exclude the
  start ($-1.47$\,dB). Seeds: 18858188808230, 2 and 3 at $\sigma_0 = 0.5$;
  11, 12 and 13 for the other CMA-ES runs.}
  \label{tab:baselines}
  \footnotesize
  \setlength{\tabcolsep}{3pt}
  \begin{tabular}{@{}l l r r r r r@{}}
    \toprule
    Optimizer, setting & Forward & Used & \multicolumn{3}{c}{Solves to} & Best \\
    \cmidrule(lr){4-6}
    & & & $-9$\,dB & $-9.5$\,dB & $-10$\,dB & (dB) \\
    \midrule
    L-BFGS, adjoint & geode-fem & $125$ & $30$ & $35$ & $60$ & $-10.04$ \\
    L-BFGS, finite differences & geode-fem & $2{,}535$ & $1{,}170$ & $1{,}365$ & $2{,}340$ & $-10.04$ \\
    \midrule
    sep-CMA-ES, $\sigma_0 = 0.5$ & geode-fem & $535$/$785$/$690$ & $405$/$385$/$395$ & $520$/$385$/$395$ & --/$695$/-- & $-9.56$/$-10.02$/$-9.77$ \\
    \quad same seeds & Palace & $860$/$780$/$660$ & $405$/$385$/$395$ & $805$/$385$/$570$ & --/--/-- & $-9.58$/$-9.89$/$-9.90$ \\
    sep-CMA-ES, $\sigma_0 = 1.0$ & geode-fem & $460$ & $380$ & -- & -- & $-9.29$ \\
    sep-CMA-ES, $\sigma_0 = 1.5$ & geode-fem & $95$/$100$/$105$ & -- & -- & -- & $-6.29$/$-1.48$/$-1.30$ \\
    \quad 1000 evaluations & geode-fem & $1{,}865$ & -- & -- & -- & $-8.81$ \\
    CMA-ES, $\sigma_0 = 0.5$ & geode-fem & $695$ & $230$ & $690$ & -- & $-9.82$ \\
    CMA-ES, $\sigma_0 = 1.5$ & geode-fem & $80$ & -- & -- & -- & $-3.51$ \\
    NEWUOA, $n + 2$ points & geode-fem & $945$ & $380$ & $400$ & $650$ & $-10.06$ \\
    NEWUOA, $2n + 1$ points & Palace & $945$ & $415$ & $420$ & -- & $-9.95$ \\
    \bottomrule
  \end{tabular}
\end{table}

\begin{figure}[t]
  \centering
  \anvilfig{fig4_optimizers.pdf}{fig4_optimizers.py}
  \caption{Reading solves to target off the optimization histories.
  (a) Best worst-of-band $|S_{11}|$ so far on the L1 optimization
  model against EM solves used, with 38 parameters, for the adjoint and
  finite-difference L-BFGS and the gradient-free runs of
  Table~\ref{tab:baselines} (the Palace CMA-ES runs, which share seeds with
  the dashed geode-fem runs, the $\sigma_0 = 1.0$ run and the short
  $\sigma_0 = 1.5$ runs are omitted). The horizontal lines are the target
  levels of Table~\ref{tab:baselines}; a run's solves to target is where its
  curve first reaches a line. The open circles mark this for $-9.5$\,dB: the
  adjoint at $35$, the fastest gradient-free run (separable CMA-ES, seed 2) at
  $385$ and the finite-difference control at $1{,}365$. (b) Solves to
  $-9.5$\,dB (solid) and $-9$\,dB (dashed) against
  the number of parameters, for the adjoint and for the best of the 5 to 11
  gradient-free runs at each $n$; the vertical bars and their labels give the
  ratio at $-9.5$\,dB. Data: \texttt{tuned/},
  \texttt{designc/c4/}, \texttt{v6/x4/}, \texttt{v6/palace\_loop/}.}
  \label{fig:optimizers}
\end{figure}

\paragraph{Results with 38 parameters.}
To reach $-9.5$\,dB the adjoint uses $35$ solves (Table~\ref{tab:baselines},
Figure~\ref{fig:optimizers}a). The fastest gradient-free run, a separable
CMA-ES seed at $\sigma_0 = 0.5$, uses $385$ ($11\times$), NEWUOA with $n + 2$
points $400$ ($11.4\times$), and no run fewer. These ratios are against
the fastest runs that reached the level; the $\sigma_0 = 1.5$ runs, censored
at $80$ to $105$ solves, bound the $-9.5$\,dB ratio only from below, by $2.3$
to $3.0\times$. The separable CMA-ES is not cross-checked against a reference
implementation, but NEWUOA, from PRIMA, gives $11.4\times$ on its own. To $-9$\,dB the ratio to the fastest run to reach it is
$7.7\times$: full CMA-ES at $\sigma_0 = 0.5$ needs $230$ solves against $30$,
in a single seed, and then $690$ to reach $-9.5$\,dB. A larger step does not help:
at $\sigma_0 = 1.5$, $89$ to $92\%$ of the samples fail the volume guard, a
190-evaluation run makes only 16 to 21 feasible evaluations, and the seed
given 1000 evaluations ($1{,}865$ solves) ends at $-8.81$\,dB, so the guard
bounds the useful step. The smaller NEWUOA interpolation set shortens the
start-up but not the time to $-9.5$\,dB ($400$ against $420$ solves); that
run also reaches $-10.06$\,dB, the best level of any 38-parameter run. The
$\sigma_0 = 0.5$ CMA-ES and NEWUOA runs had $4$ to $8$ times the adjoint's
solves, so the budgets favour them.

At $-10$\,dB the comparison is weak. Two gradient-free runs reach it, at
$650$ and $695$ solves against the adjoint's $60$, and the rest are
censored. The forward models agree to $0.004$\,dB on the 41 designs
both evaluated (the first CMA-ES generation of each seed, 42 samples with
one rejected by the guard in both) but differ by up to about $0.25$\,dB at
optimized designs (Palace reads the adjoint's final design at $-9.80$\,dB
against geode-fem's $-10.04$\,dB), the size of the gap between reaching and
missing $-10$\,dB, so we derive no ratio from the $-10$\,dB column.

\paragraph{The finite-difference control.}
With a forward-difference gradient, L-BFGS takes the adjoint's path. The two
gradients agree to $4.6\times10^{-7}$ (relative) at the start and
$5.2\times10^{-4}$ at the final design, the worst-of-band value at each of
the 12 steps matches the adjoint run's to $1.7\times10^{-5}$\,dB, and both
reach $-10.04$\,dB at step 12. A step costs $195$ solves against the
adjoint's $5$, so the ratio is $39\times$ at every level, and the $1{,}365$
solves to $-9.5$\,dB are $2.6$ to $3.55$ times the $385$--$520$ of separable
CMA-ES at $\sigma_0 = 0.5$ and more than every gradient-free run that
reached the level. With the optimizer, basis, guard and iterates fixed, only
the cost of the gradient changes, so this ablation credits the solve savings
to the adjoint; here a plain forward-difference gradient on 38 parameters is
worth less than none. Broyden-type updates reduce finite-difference
sensitivity cost in antenna optimization \citep{koziel2020expedited}, and
trust-region microwave design has been run with adjoint and with
finite-difference sensitivities \citep{koziel2012robust}; we did not try a
reduced-cost finite-difference scheme. The design itself comes from the optimizer and the
parametrization, which give the same iterates with either gradient.

\paragraph{Scaling with the number of parameters.}
From the same cold start with 4 and 14 parameters, against the best of three
separable CMA-ES seeds at $\sigma_0 = 1.5$ and NEWUOA with $n + 2$ and
$2n + 1$ points, the ratio to the fastest run to reach each level is
$1.1\times$ at $-9$\,dB and $2.4\times$ at $-9.5$\,dB with 4 parameters,
$3.3\times$ and $2.9\times$ with 14, and $7.7\times$ and $11\times$ with 38 (Figure~\ref{fig:optimizers}b). With 4 the adjoint stops on the guard
at $-9.85$\,dB and three gradient-free runs end lower ($-9.92$ to
$-9.97$\,dB); with 14 it reaches $-10.00$\,dB in $60$ solves and the best
gradient-free run ends at $-9.99$\,dB. The growth from 4 to 14 parameters
comes from the adjoint getting faster with the richer basis ($85$, $40$ and
$35$ solves to $-9.5$\,dB at 4, 14 and 38), while the best gradient-free
count is not monotone ($205$, $115$, $385$); only the step to 38 comes from
gradient-free search slowing. Three problem sizes fix the direction of the
trend but not its rate: from 14 to 38 parameters the best gradient-free count
to $-9.5$\,dB grows $3.3$-fold for $2.7$ times the parameters while the
adjoint's falls slightly, but the 4-to-14 step moves the other way for
gradient-free search, and we fit no law to three points from one problem. With 4 parameters gradient-free search is
faster to $-5$ and $-8$\,dB ($5$ and $45$ against $30$ and $55$ solves), and
the crossover is near $-9$\,dB. The adjoint's final level is set by step
rejection at the volume guard: L-BFGS with a reject-the-step guard stops at
the edge of the feasible set, while CMA-ES and NEWUOA rank or penalize
infeasible samples and keep moving along it, which is why three
gradient-free runs end lower with 4 parameters and why the $-10$\,dB column
is weak. A projected or penalized gradient step was not tried. At 4 and 14 parameters separable
CMA-ES ran at $\sigma_0 = 1.5$ only, the worst setting at 38, which may
favour the adjoint on the CMA-ES side, but NEWUOA, which has no $\sigma_0$,
gives the best gradient-free count at both. Given the 38-parameter chain's final band
shift and feed re-base in advance (Section~\ref{sec:procedure}), the
four-parameter rectangle does reach $-10.33$\,dB at \PL{}, against
$-10.78$\,dB for the 38-parameter design.

\paragraph{The forward solver and the wall-clock.}
The paired Palace runs reach $-9$\,dB at the same counts as their geode-fem
twins, so the result does not depend on geode-fem's forward model, though
Palace's band evaluation is the cheaper one ($64.4$ against $72.1$\,s on an
r6i.4xlarge). To $-9.5$\,dB the Palace-driven runs took a measured
$84$--$172$\,min (CMA-ES) and $90$\,min (NEWUOA) against an estimated
$11.1$\,min for the adjoint (its solves times the measured per-solve cost). These
are serial times on one host; CMA-ES would spread across hosts, while L-BFGS
spreads only over the five band frequencies, so solves to target is the
hardware-independent measure.

\section{Design Procedure}
\label{sec:procedure}

\paragraph{Frequency space mapping.}
A design optimized on L1 resonates higher on a finer model. We model the
fine response as the coarse response shifted in frequency,
$S_{\mathrm{fine}}(f) \approx S_{\mathrm{coarse}}(f - \delta)$, fit $\delta$
by least squares over the fine model's points, and re-optimize on the coarse
model for the band shifted down by $\delta$. For the cold L1 optimum of
Section~\ref{sec:baselines} the fitted L1-to-\PL{} shift is $47.85$\,MHz
(complex misfit $0.23$--$0.38$ unshifted, $0.018$--$0.025$ shifted); without
it, the cold L1 optimum ($-10.04$\,dB on L1) gives $-4.76$\,dB at \PL.
The design is built by a chain of such descents, each warm-started from the
previous design with the shift and, where needed, the mesh updated; we call
this the space-mapping chain. Its last descents give the three designs of
Table~\ref{tab:chain}.

\paragraph{Feed re-basing.}
Moving the feed toward the $+y$ edge compresses the mesh between them, and
the cold runs stop there at the volume guard. We then regenerate the mesh
with the port at its moved position, so that the moved feed becomes the
unmorphed base of a new mesh (re-basing), and carry the design over by
shifting the two feed coefficients. After re-basing, the optimizer widens the patch
until the guard binds in the air $5.5$\,mm outside a $\pm s$ edge, which is
a non-radiating edge.

\begin{table}[t]
  \centering
  \caption{The designs. $\delta$ is the band shift of the final descent.
  ``Solves'' counts the driven solves of that descent as recorded,
  including its gradient checks, a different count from
  Table~\ref{tab:baselines}'s. ``Fine model'' is the model the design was
  centred for.}
  \label{tab:chain}
  \small
  \begin{tabular}{@{}l l c c c l@{}}
    \toprule
    Design & Feed & $\delta$ (MHz) & Steps & Solves & Fine model \\
    \midrule
    A & $4.27$\,mm strip (free width) & $44.94$ & $9$ & $125$ & \PL \\
    B & $1.6$\,mm probe (pinned) & $44.94$ & $10$ & $130$ & \PL \\
    C (final) & $1.6$\,mm probe (pinned) & $57.0$ & $10$ & $125$ & refined rim, Palace $p = 3$ \\
    \bottomrule
  \end{tabular}
\end{table}

\paragraph{Three designs.}
Table~\ref{tab:chain} lists the designs. Design~A ends the space-mapping
chain centred for \PL{} with the feed width free, which widened the port
into a $4.27$\,mm strip. At \PL{} it reaches $-10.78$\,dB with a
$94.9$\,MHz band, $2.3938$--$2.4887$\,GHz. At Palace $p = 3$ its band moves
up to $2.4029$--$2.4975$\,GHz, the lower edge misses $2.400$\,GHz by
$2.9$\,MHz, and the worst in-band value is $-9.56$\,dB, so a design centred
on \PL{} misses once the rim is resolved. Design~B re-runs A's last step
with the probe width pinned at $1.6$\,mm and reaches $-10.62$\,dB at \PL{}
($-10.65$\,dB with a $1.27$\,mm probe), so the wide strip was not what
carried the bandwidth.

Design~C re-centres B for the refined model with a band shift of
$57.0$\,MHz. The shift adds three terms: $44.94$\,MHz from L1 to \PL,
$8.7$\,MHz for design~A's $p = 2$ to $p = 3$ dip shift in Palace, and
$3.4$\,MHz further down because every refinement so far has moved the band
up. A parallel run with the first two terms only, $53.6$\,MHz, was set
aside by a prediction from the fine models. For each run we added the
initial patch's step from the $0.2$\,mm-rim mesh to Palace $p = 3$, about
$9$\,MHz, to its lower edge on that mesh; this predicts a lower-edge margin
at $p = 3$ of about $2.6$\,MHz for the $53.6$\,MHz run and about $6$\,MHz for
C. Because Palace $p = 3$ informed both the shift and this choice, the
geode-fem mesh that was not used in it carries the independent check in
Section~\ref{sec:results}. The descent to C ran 10 steps and 125 solves, from
$-8.73$ to $-10.64$\,dB on its shifted L1 band, and stopped on the guard
with the gradient norm down from $0.143$ to $2.2\times10^{-3}$; the gradient
matched finite differences to $1.2\times10^{-6}$. As built, C is
$29.00$\,mm long and $49.72$\,mm wide on the centre line, with a barrel
outline whose $\pm s$ edges bow out most at mid-length, and its $1.6$\,mm
probe sits $5.14$\,mm from the $+y$ edge.

\section{Results}
\label{sec:results}

\begin{figure}[t]
  \centering
  \anvilfig{fig5_response.pdf}{fig5_response.py}
  \caption{(a) Return loss of design~C. Palace at $p = 3$ and $p = 2$ on the
  design's L1 mesh and at $p = 2$ on its $0.2$\,mm-rim mesh (1\,MHz grids);
  geode-fem at $p = 1$ on the combined $0.1$\,mm-rim and $0.3$\,mm-slab mesh
  (13 points) and on the $0.2$\,mm-rim mesh; and geode-fem at $p = 1$ on L1,
  the optimization model, where the design sits $57$\,MHz low by
  construction. Triangles mark the dips of Palace $p = 3$ and of L1; the
  dashed line is the $-10$\,dB target and the shading the ISM band, as in
  Figure~\ref{fig:reading}. (b) $|S_{11}|$ at $2.400$\,GHz if a whole Palace response moves up
  by a further $d$, read from the computed response at $2.400 - d$\,GHz; the
  dots mark the lower-edge margins, $6.5$ and $1.9$\,MHz. The shaded range
  is the rise of the lower edge with a $60$ or $80$\,mm air margin
  (geode-fem, $p = 1$, L1); the budget of Table~\ref{tab:budget} uses the
  dip's $4.3 \pm 1.3$\,MHz. Data: \texttt{recentre/},
  \texttt{designc/c3/}, \texttt{tuned/tuned\_summary.toml}.}
  \label{fig:response}
\end{figure}

\begin{table}[t]
  \centering
  \caption{Design~C on each model, in order of refinement, with design~A on
  Palace $p = 3$ for comparison. ``Worst'' is over the nine in-band points
  ($2.400$ to $2.4835$\,GHz in $10.4375$\,MHz steps). Band crossings are
  interpolated linearly in dB. Margins are $2.400$\,GHz minus the lower
  crossing and the upper crossing minus $2.4835$\,GHz. Design~C's meshes are
  the re-based fixtures with the port at its moved position, so their counts
  differ slightly from Table~\ref{tab:x1}'s; Palace counts $80{,}200$ more
  unknowns at $p = 2$ on the same mesh because it duplicates the vertices on
  interior PEC and port surfaces. Model-dependent values.}
  \label{tab:verify}
  \small
  \setlength{\tabcolsep}{4pt}
  \begin{tabular}{@{}l r c c c c@{}}
    \toprule
    Model & DOFs & Worst (dB) & $-10$\,dB band (GHz) & Coverage & Margins (MHz) \\
    \midrule
    geode-fem $p = 1$, L1 & $150{,}541$ & $-4.94$ & $2.3393$--$2.4304$ & $36\%$ & $60.7$ / $-53.1$ \\
    geode-fem $p = 2$, L1 & $804{,}336$ & $-9.08$ & not swept & not swept & not swept \\
    Palace $p = 2$, L1 & $884{,}536$ & $-9.08$ & $2.3847$--$2.4770$ & $92.2\%$ & $15.3$ / $-6.5$ \\
    geode-fem $p = 1$, rim $0.2$ & $431{,}528$ & $-9.23$ & $2.3849$--$2.4782$ & $93.6\%$ & $15.1$ / $-5.3$ \\
    geode-fem $p = 1$, rim $0.1$ + slab & $1{,}738{,}279$ & $-10.73$ & $2.3934$--$2.4881$ & $100\%$ & $6.6$ / $4.6$ \\
    Palace $p = 3$, L1 & $2{,}503{,}077$ & $-10.46$ & $2.3935$--$2.4864$ & $100\%$ & $6.5$ / $2.9$ \\
    Palace $p = 2$, rim $0.2$ & $2{,}465{,}984$ & $-10.32$ & $2.3981$--$2.4912$ & $100\%$ & $1.9$ / $7.7$ \\
    \midrule
    Design A, Palace $p = 3$ & $2{,}729{,}331$ & $-9.56$ & $2.4029$--$2.4975$ & $96.5\%$ & $-2.9$ / $14.0$ \\
    \bottomrule
  \end{tabular}
\end{table}

\paragraph{On the finest models.}
On Palace at $p = 3$ on L1, design~C has a worst in-band $|S_{11}|$ of
$-10.46$\,dB, at $2.4835$\,GHz, and a contiguous $-10$\,dB band from
$2.3935$ to $2.4864$\,GHz, $93.0$\,MHz wide, which covers the $83.5$\,MHz ISM
band with $6.5$\,MHz to spare at $2.400$\,GHz and $2.9$\,MHz at $2.4835$\,GHz
(Table~\ref{tab:verify}, Figure~\ref{fig:response}a). geode-fem's finest
mesh, the combined $0.1$\,mm rim and $0.3$\,mm slab ($1.74$M unknowns),
gives $-10.73$\,dB and a $94.7$\,MHz band, $2.3934$--$2.4881$\,GHz. The two
lower edges agree to within the sweep resolution (the interpolated values
differ by $0.02$\,MHz); the upper edges differ by $1.7$\,MHz, partly
because geode-fem's upper crossing is interpolated across a $5$\,MHz gap. Palace at $p = 2$ on the $0.2$\,mm-rim mesh ($2.47$M unknowns), the
most refined model computed for the design, gives $-10.32$\,dB and a band
of $2.3981$--$2.4912$\,GHz: its lower edge sits $4.6$\,MHz above Palace
$p = 3$ on L1, which leaves $1.9$\,MHz of lower-edge margin. The two models
are different refinements of comparable size (one more order on L1, and the
rim at $p = 2$), so the $4.6$\,MHz shows that $p = 3$ on L1 does not yet
resolve the rim. From $p = 2$ to $p = 3$ on L1 the design's dip rises
$9.1$\,MHz and its lower band edge $8.7$\,MHz; the initial patch's dip rises
$8.9$\,MHz over the same step, so the design shifts like the patch.

On the refined meshes the air is moved on cylinders (the default morph
inverts elements on the combined mesh); on the $0.2$\,mm-rim mesh the two
air motions agree to $0.09$\,MHz at the lower edge.

\begin{table}[t]
  \centering
  \caption{What moves design~C's lower $-10$\,dB edge, against its margin on
  Palace $p = 3$, L1. A negative entry uses up margin.}
  \label{tab:budget}
  \small
  \begin{tabular}{@{}>{\raggedright\arraybackslash}p{0.43\linewidth} r >{\raggedright\arraybackslash}p{0.33\linewidth}@{}}
    \toprule
    Term & MHz & Tier, source \\
    \midrule
    Lower-edge margin, Palace $p = 3$, L1 & $+6.5$ & model-dependent, Table~\ref{tab:verify} \\
    Rim refinement at $p = 2$ (Palace, $0.2$\,mm rim) & $-4.6$ & model-dependent, Table~\ref{tab:verify} \\
    Air margin, $40 \to 60$ / $80$\,mm (dip; edges scatter by $1.3$) & $-4.3$ & model-dependent at $p = 1$ on L1, level by $60$\,mm; transfer estimated \\
    Further refinement above the $0.2$\,mm-rim model & not computed & initial patch: $6.4$--$13.5$ above Palace $p = 3$, estimated \\
    Coaxial feed in place of the lumped port & not quantified & \\
    Copper loss in place of PEC & not quantified & \\
    \bottomrule
  \end{tabular}
\end{table}

\paragraph{The lower-edge budget.}
Table~\ref{tab:budget} collects what moves the lower edge. We take the
air-margin term as the dip's $4.3$\,MHz and carry the $1.3$\,MHz edge scatter
between re-meshed air boxes (Section~\ref{sec:accuracy}) as the uncertainty
of the measured terms; it is the scatter of the air re-meshes only, applied
once to the sum. Added, they take more than the margin:
$6.5 - 4.6 - 4.3 = -2.4 \pm 1.3$\,MHz, a lower edge near $2.402$\,GHz, and on
the $0.2$\,mm-rim response moved up by $4.3$\,MHz, $2.400$\,GHz sits at
$-9.63$\,dB (Figure~\ref{fig:response}b). The sum assumes that the terms add,
and the air-margin term was measured on the coarse model. The rim step
($4.6$\,MHz) is below both drift estimates of the initial patch above Palace
$p = 3$ (order-1 Richardson $6.4$\,MHz, rim fit $13.5$\,MHz), so it need not be
the last of them; Palace $p = 3$ on the $0.2$\,mm-rim mesh was not run. If the
estimates transfer and include the rim step, up to $13.5 - 4.6 = 8.9$\,MHz
more remains. The estimated miss is therefore $2.4 - 1.3 = 1.1$\,MHz to
$2.4 + 1.3 + 8.9 = 12.6$\,MHz, which we state as $1$ to $13$\,MHz. A further
re-centring cannot cover the whole range: the $93.0$\,MHz band is $9.5$\,MHz
wider than the ISM band, less than the $11.5$\,MHz spread of the estimate, so
the bandwidth, which the guard and the board limit, binds.

\paragraph{What is claimed.}
Measured on stated discrete models (model-dependent): design~C's $-10$\,dB
band covers $2.400$--$2.4835$\,GHz on the three most refined models computed,
Palace $p = 3$ on L1, geode-fem's combined $0.1$\,mm-rim and $0.3$\,mm-slab
mesh and Palace $p = 2$ on the $0.2$\,mm-rim mesh, with $6.5$, $6.6$ and
$1.9$\,MHz to spare at $2.400$\,GHz. Estimated: a larger air margin and
further refinement each move the band up, and the budget above puts the
mesh-limit lower edge $1$ to $13$\,MHz above $2.400$\,GHz. The claim is
therefore ISM coverage on those three discrete models. We do not claim ISM
coverage for the mesh-limit model, whose lower edge we estimate misses the
band, or for a built antenna, whose coaxial feed and copper loss are not
modelled (Section~\ref{sec:limits}).

\paragraph{Where the descent stopped.}
The descents to A, B and C stop on the volume guard. For C the gradient norm fell 65-fold first, so C may be
near a stationary point of the guarded problem (for the rectangle it grew,
$0.023$ to $0.074$), and the widened patch comes within about $6$\,mm of
the $62$\,mm board edge. Remeshing at the widened outline would relieve the
guard; it was not run.

\paragraph{Radiation efficiency and gain.}
For design~C we computed the accepted port power, the Poynting flux through
a closed surface in the air box, the FR-4 loss and the broadside gain by a
near-to-far-field transform at $p = 1$, on L1 at the band it was optimized
on and on the $0.2$\,mm-rim mesh across the ISM band. The power balance
closes to $0.8\%$ on L1 and $1.6\%$ on the rim mesh. The radiation
efficiency is $0.61$--$0.63$, the dielectric loss fraction $0.37$, and the
broadside gain $4.45$--$4.66$\,dBi with a front-to-back ratio of
$13$--$15$\,dB. With each model's own match the realized broadside gain is
$4.06$--$4.50$\,dBi on L1 over its optimization band and $4.11$--$4.56$\,dBi
on the rim mesh over the ISM band (model-dependent). With perfectly
conducting (PEC) conductors this efficiency counts radiation against dielectric loss only;
copper loss would lower it and the gain.

\paragraph{Bandwidth against a closed form.}
On the voltage-standing-wave-ratio-2 criterion, Palace gives design~C a
bandwidth of $99.0$\,MHz at $p = 3$ (design~A: $100.3$\,MHz at $p = 2$ and
$101.0$\,MHz at $p = 3$). The closed-form formulas of
\citet{jackson1991simple}, evaluated for design~A's $49.75\times28.80$\,mm
rectangle (C's outline is $49.72\times29.00$\,mm), give $70.4$\,MHz, and
$67.9$\,MHz at the area-equivalent width of $45.4$\,mm. Since $\tan\delta$ is
the only loss, design~C's dielectric loss fraction ($0.367$--$0.372$) gives the
model inequality $Q_t \ge 0.3667/\tan\delta = 18.3$ (not a bound in the sense of
Section~\ref{sec:intro}), against the closed form's $Q_t = 24.5$ and radiation
$Q$ of $48.0$ for design~A's rectangle. We read this, as an interpretation, as
the model patch radiating more strongly than the infinite-ground closed form,
perhaps through its finite $62\times53$\,mm ground, which we did not test.

\paragraph{Cost.}
The final descent took $26$\,min on a 32-vCPU host, the finite-difference
control $4.7$\,h on a 16-vCPU host ($1{,}164$ to $1{,}211$\,s per gradient with 8
finite-difference jobs in parallel), the Palace $p = 3$ sweep of design~C $3.5$\,h on 16 MPI ranks and
the $0.2$\,mm-rim sweep $1.5$\,h. geode-fem's direct solver needed about
$8$\,min per frequency at $204$\,GiB on the combined mesh.

\section{Limitations and Future Work}
\label{sec:limits}

\paragraph{Feed, losses, refinement and scope.}
The probe is a uniform lumped port of the pin's width across the substrate,
not a coaxial probe. A coaxial feed adds the pin's series inductance and a
coax-to-substrate transition. Loading the gap with the lumped port alone
moves design~C's mode by $23$\,MHz, so a coax-fed model could move the
match by megahertz, as much as the margins; Palace supports wave ports, and
a coax-fed model was not run. Conductors are PEC, and the design has not
been fabricated or measured.
The drift above the $0.2$\,mm-rim $p = 2$ model is estimated, not computed.
A $p = 2$ shape adjoint, or optimizing on a rim-refined mesh, would remove
the coarse-to-fine shift that space mapping corrects here, and re-centring
with the air-margin term included is the next step. The optimizer
comparison is one problem from one start on the coarse model, with one
full-CMA-ES seed per step size and one NEWUOA trust radius; a
surrogate-assisted or Bayesian optimizer may narrow the gap, and whether the
adjoint helps at the refined level is not tested.

\section{How geode-fem Was Designed and Built}
\label{sec:built}

geode-fem implements the whiteroom operator specification of Crutcher
Dunnavant \citep{dunnavant2026whiteroom}, whose four layers restate Palace's
cited source as pure functions, fused primitives, tensor operations and a
small graph-evaluation calculus. Dunnavant also gave design input on the
solver's operator layer and contributed code: the migration to Rust 2024, the
Metal (Apple GPU) backend, the present layout of the core crate and the
examples, and the geode-util crate. The rest of the code, the experiments and
the text of this paper were produced by AI models working as role-separated
agents: Anthropic's Claude (Opus, Fable and Sonnet models, through Claude
Code) and OpenAI's GPT Sol and GPT Astra models (via Codex). The agents were
coordinated by Loom \citep{loom}, an issue, worktree and pull-request
workflow, and by Squad, a message room the agents share. Nearly every
change was reviewed by a separate agent session. The agent reviews of the
experiments behind this paper recomputed the reported numbers from the
archived records, and the script that checks this paper's text against
those records was mutation-tested. The human author directed the research,
made every scoping and claim decision, and takes responsibility for the
content.

\section{Reproducibility}
\label{sec:repro}

Every result in this paper is read from an archived, scripted
artifact.\footnote{The geode-fem source is public at
\url{https://github.com/rjwalters/geode-fem}.}
% Archival DOI: add at deposit.
Table~\ref{tab:artifacts} lists them, relative to
\path{benchmarks/patch_antenna_conformal/}. Each run record holds the exact
command and the SHA-256 of its mesh; the meshes are regenerated from YAML
recipes in \path{reference/gmsh/} with gmsh 4.15.2, and runs on re-based
fixtures need the \texttt{--feed-support-legacy} flag. Many records carry no
generating commit (they ran outside a git checkout); the code is merged at
geode-fem commit \texttt{99646f06}, and the setting sweep and air-margin
runs used \texttt{dcd8df60}. Palace is commit \texttt{fba6a5b}, built from
\path{reference/palace/docker/Dockerfile}; PRIMA was built from libprima
commit \texttt{43863c69} (package \texttt{0.7.3.dev1801+g43863c69c}), and
pycma is 4.5.0. geode-fem ran on AWS Graviton4 hosts (r8g.4xlarge and
r8g.8xlarge), Palace and the cost measurements on r6i.4xlarge. The figure
scripts in \path{figures/src/} read the artifacts, \path{provenance.md} maps
every number to its source field, and \path{verify_evidence.py} checks the
text against them.

\begin{table}[t]
  \centering
  \caption{Artifacts behind each result, relative to
  \texttt{benchmarks/patch\_antenna\_conformal/} except where a path begins
  with \texttt{benchmarks/}, \texttt{reference/} or \texttt{crates/} (repository
  root).}
  \label{tab:artifacts}
  \footnotesize
  \begin{tabular}{@{}>{\raggedright\arraybackslash}p{0.30\linewidth}>{\raggedright\arraybackslash}p{0.64\linewidth}@{}}
    \toprule
    Result & Artifact \\
    \midrule
    Fixture, refinement, Palace cross-check, curvature (Section~\ref{sec:accuracy}) & \path{baseline_2g4_L1.toml}, \path{convergence.toml}, \path{upml_ladder_L1.toml}; \path{v6/v6_summary.toml}, \path{v6/x1/}, \path{v6/x2/} \\
    Baselines and scaling (Section~\ref{sec:baselines}) & \path{tuned/tuned_summary.toml}, \path{tuned/}; \path{v6/palace_loop/palace_loop_summary.toml}, \path{v6/x4/}; \path{copper/L1_copper_38.toml}; \path{designc/c4/} \\
    Designs A to C (Sections~\ref{sec:procedure}, \ref{sec:results}) & \path{copper/g2_summary.toml}, \path{v6/x3/}, \path{v6/x5/}; \path{recentre/recentre_summary.toml}, \path{recentre/palace/}; \path{designc/designc_summary.toml}, \path{designc/c1/}--\path{c3/}; \path{tuned/b5/} \\
    FDTD cross-check (Section~\ref{sec:accuracy}) & \path{reference/meep/docker/DIAGNOSIS.md}, \path{reference/meep/docker/HARNESS_V2.md}; \path{benchmarks/fdtd_density_baseline/meep_runs/v2/ladder.json} \\
    Generator and tests & \path{crates/geode-core/examples/patch_conformal_nested_dof.rs}; \path{crates/geode-core/tests/patch_conformal_*.rs} \\
    \bottomrule
  \end{tabular}
\end{table}

\section{Conclusion}
\label{sec:conclusion}

On the copper design of a bent 2.4\,GHz patch, a finite-element shape
adjoint reaches a worst-of-band $|S_{11}|$ of $-9.5$\,dB on the coarse
optimization model with $11$ times fewer EM solves ($8.3$ times less geode-fem
wall time) than the fastest gradient-free run to reach that level, and
$-9$\,dB with $7.7$ times fewer than the fastest run to reach $-9$\,dB. The
ratio at $-9.5$\,dB is $2.4\times$, $2.9\times$ and $11\times$ with 4, 14
and 38 parameters: it rises with their number, first because the adjoint gets
faster with a richer basis and then because gradient-free search slows, and
three sizes of one problem fix no rate. A finite-difference gradient on the
same path costs more than gradient-free search, so the savings are the
adjoint's. The optimum is no better, since the guard stops this L-BFGS where
gradient-free search reaches the same level given enough solves. The design,
re-centred using Palace at third order and a rim-refined mesh, with a
$1.6$\,mm probe on a standard board, covers the ISM band at $-10$\,dB on the
three most refined models computed, with $1.9$\,MHz to spare on the most
refined; that coverage is what we claim. A larger air margin moves the
resonance another $4.3$\,MHz and further refinement moves it up again, so we
estimate that the mesh-limit design misses $2.400$\,GHz by $1$ to $13$\,MHz.

\paragraph{Acknowledgments.}
We thank Crutcher Dunnavant for the whiteroom specification, for his
design input to the solver's operator layer and for his contributions to
the codebase. The AI-orchestrated development of geode-fem and of this
paper is described in Section~\ref{sec:built}.

\bibliographystyle{plainnat}
\bibliography{refs}

\end{document}
